\clearpage\section{Moodle integration and security}\label{sec:security}
\subsection{Identity linking and course authority}
The integration boundary consists of Moodle Web Services\citep{moodlews} and the \code{local_lms_bridge} plugin. The bridge provides identity, capability and publication contracts.

Local application identity and Moodle identity are independent. Registration does not prove control of a Moodle account. A link request produces an opaque request ID; a matching Moodle user must sign in and confirm on the Moodle side. The application then stores the confirmed Moodle user ID. A generic pending response prevents the link-request response from revealing whether the entered email matches an active account. Missing and ambiguous emails receive pending requests with the same externally visible shape; only legitimate confirmation may reveal identity.

A confirmed identity does not automatically grant access to every Moodle course. The bridge queries course authoring capabilities and verifies the requested actor in the course context. The author ID is propagated on creation calls. The bridge also checks relevant activity/question capabilities in its external handlers. The application administrator is not treated as a Moodle administrator solely because the local admin flag is true. Rights can be revoked externally, so publication verifies them immediately before activity writes.
\subsection{Publication mapping and sequence}
The Book receives navigable lesson material, including the LO and usable generated content. Practice is published as a quiz with immediate feedback; the assessment Quiz uses deferred feedback. Question stems, choices and explanations are exported from the live saved question data. The target section is chosen by the author. Moodle's Book and Quiz remain the actual learner-facing activities\citep{moodlebook,moodlequiz}.
Figure~\ref{fig:moodle} shows publication sequence and permission checks. Book and subsequent quiz writes are separate remote operations.
\diagram{moodle-sequence}{Publication sequence and permission checks. Book and subsequent quiz writes are separate remote operations.}{moodle}
The application sanitizes HTML through an allowlist before transmission and returns created activity IDs. Book, Practice and Quiz are separate remote operations. If a later operation fails, earlier activities can remain. The publication design records external IDs and per-step state, expose partial success and support safe resumption during recovery.
\subsection{LTI 1.3 launch}\label{sec:lti}
The advisor asked for the Moodle-to-LTI-tool model. GENAILMS now acts as an LTI 1.3 tool: an instructor opens an External tool activity in a Moodle course and arrives signed in, in the matching course, without a password. The bridge plugin is kept for publishing Book and Quiz activities, because a standard LTI tool cannot create them.
\begin{enumerate}[nosep]
\item Moodle calls \code{/api/lti/login}. The tool stores a \code{state} and \code{nonce} in PostgreSQL with a 10-minute expiry and redirects to Moodle's authorization endpoint.
\item Moodle posts a signed \code{id_token} to \code{/api/lti/launch}. The tool consumes the state, so it is usable once even when the launch is rejected, and verifies the RS256 signature with Moodle's key set, which is cached and fetched again when an unknown key id appears.
\item It checks issuer, audience (and \code{azp} when there are several audiences), expiry and issue time with 60 s of clock skew, nonce, a registered deployment, version \code{1.3.0}, message type \code{LtiResourceLinkRequest} and the presence of a course context.
\item It finds or creates the user by the pair of platform and \code{sub}, and the course by platform, deployment and context id, then redirects to the frontend with a one-time code in the URL fragment, which the frontend exchanges for the same tokens a password login gives.
\end{enumerate}
Email does not identify an LTI user, because a launch does not prove control of an address here. For the project's own Moodle, the LTI \code{sub} and context id are Moodle's user and course ids, guaranteed by Moodle's signature, so a launch reuses the account and course already linked through the plugin instead of creating second ones. The launch was run end to end against Moodle 5.2.3 in an isolated stack: the instructor entered and saw the course, a second launch reused the same account and course, a student was refused, and the stored Moodle ids matched. Rejection cases (wrong signature, audience, expiry, nonce or deployment, replayed state, reused code) are tested with a fake platform signing with a real RSA key.

Current limits: each course belongs to one instructor; Deep Linking and grade return are not implemented, because the tool publishes modules through the plugin and does not grade students; and the login state is not bound to the browser, so an attacker could make a victim sign in to the attacker's own tool account (login cross-site request forgery). Binding the state to a \code{SameSite=None} cookie when the tool opens in a new window would prevent it. The tool should be opened in a new window: inside an iframe, browsers partition storage, so the session would not be shared with a directly opened tab.
\subsection{Student access through LTI}\label{sec:lti-learners}
A student who opens the same External tool activity may chat with the course's material and do nothing else. Three decisions were made before the change: the lecturer does not read students' chats, a student may send 30 chat messages per day, and a student's launch is refused until the instructor has opened the tool once, because the course does not exist in the tool before that.

Until this change, ``a course I can see'' and ``a course I own'' were the same condition in every row-level security policy. Course membership separates them (Table~\ref{tab:member-access}):
\begin{itemize}[nosep]
\item A \code{course_members} table holds (course, user) pairs. A row can be inserted only by the member themself and only for a course that an LTI context opens, so membership comes only from a signed launch. The \code{courses} policy reads membership through a \code{SECURITY DEFINER} function, because PostgreSQL refuses a policy on \code{courses} that reads \code{course_members} while the policy on \code{course_members} reads \code{courses}.
\item Members get policies for \code{SELECT} only, on \code{courses}, \code{documents} and \code{chunks}. Every write still requires the owner; modules, jobs and their children stay owner-only.
\item \code{chat_sessions} gained a user column. A session is visible only to the person who had it, in a course that person can see; existing sessions were assigned to the course owner, who had them.
\item \code{users.can_author} is false for an account created by a Learner launch and set true by an Instructor launch. Creating a course requires it, and the daily chat counter applies only to accounts that cannot author.
\end{itemize}
\begin{table}[H]\centering\small\caption{What a course member can reach.}\label{tab:member-access}
\begin{tabularx}{\linewidth}{|p{4.2cm}|Y|Y|}\hline
Resource & Owner & Member through LTI\\\hline
Course, documents, chunks & Read and write & Read, in the database; the API exposes them only through chat\\\hline
Modules, jobs, publication & Read and write & None\\\hline
Course chat & Own sessions & Own sessions, 30 messages per day\\\hline
Other people's chats & None & None\\\hline
\end{tabularx}\end{table}
The API keeps the owner check on every route except chat, which accepts an owner or a member. \code{/lti/session} returns the course role, and the frontend opens a member's launch on a chat-only page; an author page reached by a student account redirects there. A student who requests a route the lecturer can reach receives 404, as for any missing resource.

The policies are tested as SQL under the runtime role: what a member reads, that it cannot change or add, and who sees which chat. Loosening the \code{documents} policy back to its previous condition makes three of these tests fail. Launch tests cover a student before the instructor (refused), a student who joins and can only chat, one membership after repeated launches, the daily limit, an author who launches as a learner and a student who later becomes an instructor. Against Moodle 5.2.3, a student before the teacher received 400; the teacher became the owner; the student then became a member, chat returned 200, the document list 404, and one membership row was stored.

A launch adds a member but never says that one has left. The worker therefore reads the course roster through the LTI Names and Role Provisioning Services 2.0. The roster address comes from the \code{namesroleservice} claim of the latest launch and must have the issuer's origin, because the access token is sent to it. The tool obtains that token with an OAuth~2 client-credentials grant, signing its assertion with its own RSA key; the public key is served at \code{/api/lti/jwks} with an RFC~7638 thumbprint as its key id. Every five minutes the worker reads, page by page, the roster of each course not checked for an hour, and removes the members who are no longer listed as active learners or instructors; a suspended enrolment counts as gone. The next request of a removed student finds no course, because membership is enforced by row-level security. Nobody is removed when the result is uncertain: a failed or refused request, a next page outside the platform, an empty roster or a course whose context has no roster address leaves membership unchanged until the next check.

Tests use a fake platform that verifies the assertion with the key set the tool publishes, and cover departure, suspension, roles the tool does not serve, several pages, one read per cycle and each case that must remove nobody. Against Moodle 5.2.3, with a key generated for the run, the worker read the roster and removed nobody while the student was enrolled, then removed exactly that student after they were unenrolled; the teacher remained the owner. A student therefore loses access within about one hour and five minutes of leaving the course. Where the platform does not offer the service or the tool has no key, the earlier limit applies: the account keeps chatting until the session of its last launch expires, at most 14 days later.
\subsection{Authentication and session security}
Passwords are stored as bcrypt hashes\citep{bcrypt}. Access tokens are JSON Web Tokens\citep{jwt} valid for 15 minutes and kept in memory by the browser; the active and admin flags are read from the database on every request, so disabling an account takes effect at once. Refresh tokens are random values stored as SHA-256 hashes, rotated on each use and sent in a Secure, HttpOnly, SameSite cookie limited to the refresh endpoint; renewal also checks the origin and a CSRF token. Logout revokes the refresh token. Rendered Markdown is sanitised and a restrictive content security policy limits cross-site scripting. Production settings refuse default secrets.
\subsection{Row-level security and ownership}
Authorisation is checked on the server and denies by default\citep{owaspaccess}, in the API and again by PostgreSQL row-level security (RLS). Table owners bypass row policies\citep{postgresrls}, so requests use a separate runtime role, and the user id is set as a transaction-local setting at the start of every transaction. Course-owned content is filtered by policies; course members have separate read-only policies (Section~\ref{sec:lti-learners}). The worker has a separate role with restricted table grants; it is a trusted service boundary, not a user session. It is exempted from row-level security by an explicit policy on each table rather than by the \code{BYPASSRLS} attribute, which the non-superuser master role of RDS cannot set (Section~\ref{sec:architecture}). No application role has \code{BYPASSRLS}, and a test fails when a new table with row-level security lacks the worker policy.

RLS filters rows but cannot hide columns, so raters read studies through views that omit the group and the source module. A rater can read and write only their own ratings. Administrative job retry goes through a narrow database function, which gives the administrator no access to course content.
\subsection{Indirect prompt injection}\label{sec:prompt-injection}
Uploaded material is untrusted input to the model: a document can contain text written as an instruction, such as ``ignore the rules above''\citep{owaspprompt}. Three controls are implemented. Every prompt that receives source passages wraps each passage in a block marked as data, not instructions, and contains a rule that instructions found inside the sources must be treated as course content and not followed. Lecturer guidance is wrapped the same way and may change emphasis and wording but not the grounding rules. The model's output is never trusted: it is validated against a schema, its citation codes must exist in the prompt, HTML sent to Moodle passes an allowlist, and the model decides no authorisation and starts no publication. Two further controls are designed but not implemented: flagging instruction-like passages at ingestion for the lecturer, and datamarking the source text\citep{spotlighting}. Delimiters and rules reduce the risk; they are not a security boundary, which is why authorisation never depends on model output.
\subsection{Threats, controls and verification targets}
Table~\ref{tab:07-moodle-security-1} maps each threat to its control and its test.
\begin{longtable}{|p{3.0cm}|p{5.8cm}|p{5.1cm}|}\caption{Security controls and their limits.}\label{tab:07-moodle-security-1}\\\hline
Threat & Control & Verification\\\hline\endfirsthead
\hline Threat & Control & Verification\\\hline\endhead
Cross-owner object access & Ownership checks, course/document filters, composite FKs and RLS; uniform missing-resource response & Direct API and runtime-role DB probes across owners and after transaction restart\\\hline
Account enumeration & Generic Moodle pending link response & Equivalent status and expiry for matching, missing and ambiguous emails; transport errors remain distinct\\\hline
Unsafe rendered content & Moodle HTML allowlist; structured model output & Markdown preview and exported activity XSS cases, including source text\\\hline
Resource exhaustion & Streamed upload limits, JSON body limit, route limits and daily generation cap for untrusted accounts & Concurrent uploads/jobs, provider quota exhaustion, CPU/RAM and fair scheduling\\\hline
Indirect prompt injection & Sources and guidance delimited as data, a rule against embedded instructions, output validation (Section~\ref{sec:prompt-injection}) & Adversarial documents\\\hline
Credential exposure & Provider keys stay at gateway; separate DB roles; secret-aware configuration fields & Production configuration, logs, rotation and backup access\\\hline
Publication misuse & Confirmed actor plus fresh Moodle capability checks & Unlinked, revoked, suspended and wrong-course scenarios\\\hline
Student overreach & Membership only from a signed launch; member policies for \code{SELECT} only; chats per person; 30 messages per day (Section~\ref{sec:lti-learners}) & Runtime-role SQL tests and Moodle 5.2.3 run done; roster checks remove a departed student within about an hour\\\hline
Blinding failure & Restricted views, own-rating policies and snapshots & Response field audit and packet-text leakage review\\\hline
\end{longtable}
JSON bodies are limited to 1 MB and uploads have their own streamed limits (NFR07). Production hides the API documentation and sets security headers. Selected source text is sent to the model provider; embedding stays on the project's own server. The optional Jev reranker would send the query and candidate chunks to TypeSafe in the United States, so it stays disabled until the university agrees.
